% HD-N26. GPT-6 Astra Pro /root/imprimitive, 21 November 2026.
% Inherited effort; numerical setting not exposed.
\section{HD-N26: all even-cohomology complexes on the positive graph supports}
\label{sec:even-cohomology-graph-complexes}

Let $C_1,\ldots,C_N$, $N=9408$, be any nodal translated
abelian-threefold arrangement of HD-N23. Consider all actual
objects $E\in D^b_{\rm coh}(A)$ with the following cohomology:
each $\mathcal O_{C_\nu}$ occurs exactly once, in an even
degree $a_\nu$, and there are no other cohomology summands.
Thus
\[
 \mathcal H^a(E)=\bigoplus_{a_\nu=a}\mathcal O_{C_\nu},
 \qquad a_\nu\in2\mathbf Z.                              \tag{N26.1}
\]
This specifies a class of actual derived objects; the
assertions below apply to every differential and every
higher gluing map of any object in this class. It includes
the actual nonzero couplings of HD-N25. Such objects are
perfect because the ambient variety is smooth, and their
whole character is
\[
                   \operatorname{ch}(E)=32\theta^3+120T.
                                                               \tag{N26.2}
\]

Every object in (N26.1) has
\[
 \operatorname{Ext}^2(E,E)\simeq
   \bigoplus_{\nu=1}^N
       \operatorname{Ext}^2(\mathcal O_{C_\nu},\mathcal O_{C_\nu}),
 \qquad \dim\operatorname{Ext}^2(E,E)=15N=141120.          \tag{N26.3}
\]
The isomorphism is the natural edge identification of the
canonical filtered-Hom construction specified below. In
particular it is compatible with every actual coherent
point action. Every faithful finite point subgroup acting
coherently on $E$ therefore has a weak degree-two trace
kernel of dimension at least $495$. This bound also holds
after any actual isogeny pullback, with every newly available
coherent point action, and after subsequent actual line
tensor and Fourier--Mukai equivalences.

The gluable objects among (N26.1) have an exact description:
\[
 \operatorname{Ext}^{<0}(E,E)=0
 \quad\Longleftrightarrow\quad
 C_\mu\cap C_\nu=\varnothing
       \text{ whenever }|a_\mu-a_\nu|\ge4.               \tag{N26.4}
\]
Both assertions concern every actual object with the stated
cohomology; no existence of a prescribed collection of
formal higher differentials is assumed.

\paragraph{1. The canonical filtered-Hom map and its source.}
Apply BBD in the abelian category of $\mathcal O_A$-modules,
which has enough injectives, and take $D^b_{\rm coh}(A)$
as the full subcategory with bounded coherent cohomology.
Use the increasing canonical filtration
$T_iE=\tau_{\le i}E$, or the decreasing filtration
$F^i=T_{-i}$. Its graded object is
\[
                  \operatorname{Gr}_F^iE
                            =\mathcal H^{-i}(E)[i].      \tag{N26.5}
\]
The precise published source is Beilinson--Bernstein--Deligne,
\emph{Faisceaux pervers}, Ast\'erisque 100 (1982),
Section 3.1.3, equation (3.1.3.4), printed p.78, together with
Proposition 3.1.6 and its proof, printed p.79. The latter says
that forgetting a canonical filtered object is an equivalence
onto the usual derived category. Thus the lift of an actual
object, every isomorphism, and their compositions are
functorial. In particular this lift retains the coherence
maps of an actual point action.

For clarity, the source construction is the filtered complex
$R\operatorname{Hom}(K,L)$: after taking a degreewise split
filtered injective model for $L$, its degree-$n$ maps in
filtration $p$ carry $F^iK$ into $F^{i+p}L$.
Its associated graded is the sum of the actual derived Hom
complexes between graded objects with filtration difference
$p$. Equation(3.1.3.4) then gives
\[
 E_1^{p,q}=
  \bigoplus_{i\in\mathbf Z}
  \operatorname{Ext}^{\,2p+q}
       (\mathcal H^i(E),\mathcal H^{i-p}(E))
  \ \Longrightarrow\ \operatorname{Ext}^{p+q}(E,E),
 \quad d_r:E_r^{p,q}\longrightarrow E_r^{p+r,q-r+1}.
                                                               \tag{N26.6}
\]
Indeed, substituting (N26.5) in the source's
$\operatorname{Hom}^{p+q}$ of graded objects gives precisely
the exponent $2p+q$ in (N26.6). The cohomological filtration
has finite length, so convergence is a finite-filtration
statement. BBD Section 3.1.1, printed p.76, gives the required
filtered injective model. The same substitution is written
explicitly in Pirozhkov's 2020 author thesis,
\emph{Admissible subcategories of del Pezzo surfaces},
Lemma 2.4.1 and proof, printed p.15; the published BBD
construction is the input here.

This proves the exact naturality needed: t-exact point
functors preserve the canonical truncations; their
isomorphisms induce maps of (N26.6), with the given
compositions. Scalar coherence cancels on conjugation of
self-Hom. No choice of an ordinary direct-sum projector on
$E$ enters this argument.

\paragraph{2. Every degree-two differential vanishes at the diagonal.}
The actual transverse Koszul maps and complement isogeny
in HD-N25 give
\[
 \begin{split}
 &\operatorname{Ext}^j(\mathcal O_{C_\mu},\mathcal O_{C_\nu})=0
                 \quad(\mu\ne\nu,\ j\ne3),\\
 &\dim\operatorname{Ext}^j(\mathcal O_{C_\nu},\mathcal O_{C_\nu})
                                         =\binom6j.
 \end{split}                                                  \tag{N26.7}
\]
The nonzero cross group in degree three is the direct sum
of its one-dimensional local Koszul groups over the actual
intersection points. Each support occurs in only one
cohomology degree. Consequently, writing $k=p+q$ for total
degree, the possible nonzero cells of (N26.6) are exactly
of the following two forms:
\[
 \begin{array}{c|c|c}
 \text{source of a summand}&p&k\\ \hline
 \text{self-Ext of one support}&0&0,1,\ldots,6\\
 \text{cross Ext of distinct supports}&p\in2\mathbf Z&3-p.
 \end{array}                                                  \tag{N26.8}
\]
The second line records possible cells, with actual dimension
zero when the corresponding supports are disjoint.

In total degree two the first page therefore has only the
cell $(p,q)=(0,2)$, and it is the sum on the right of
(N26.3). An outgoing $d_r$ from this cell has
$p=r$, total degree $3$, and ordinary cross-Ext degree
$r+3\ge4$. Because $r>0$, no same-support term occurs
there. It is zero by (N26.7). An incoming $d_r$ starts at
$p=-r$, total degree $1$, with ordinary cross-Ext degree
$1-r\le0$. Its supports lie in distinct cohomology
degrees, so it too is zero. These are zero already on
$E_1$, hence on every later page. Thus
\[
                   E_\infty^{0,2}=E_1^{0,2},\qquad
                   E_\infty^{p,2-p}=0\quad(p\ne0).       \tag{N26.9}
\]
The finite filtration on $\operatorname{Ext}^2(E,E)$ has
just this one nonzero graded quotient. Its edge
identification is (N26.3), natural under the filtered maps
of paragraph1. This proves the stated equality for every
actual higher coupling, not only a two-level cone.

\paragraph{3. Negative self-Ext and the gluable locus.}
A negative total-degree cell can only come from the second
line of (N26.8), with $p\ge4$ even and $k=3-p<0$.
Its outgoing $d_r$ has filtration $p+r$ and total degree
$4-p$, hence ordinary Ext degree $r+4\ge5$; the target
supports are distinct and that group is zero. An incoming
$d_r$ has filtration $p-r$ and total degree $2-p$.
When $p-r\ne0$, its ordinary cross-Ext degree is
$2-r\le1$, so it vanishes. When $p-r=0$, it is a
negative same-grade self-Ext and also vanishes. Thus every
negative first-page cell survives. In particular, for $k<0$,
\[
 \dim\operatorname{Ext}^k(E,E)=
 \sum_{\substack{\mu,\nu\\a_\mu-a_\nu=3-k}}
                         \#(C_\mu\cap C_\nu),           \tag{N26.10}
\]
where all intersections are the transverse finite ones of
HD-N23; disjoint pairs contribute zero. There is at most
one filtration index in each such total degree, so the
formula follows directly from the finite abutment.
This proves (N26.4), including both directions. The
two-level choices in HD-N25 have differences only $0,2$,
and hence recover its gluable conclusion.

\paragraph{4. Every coherent point action and every isogeny pullback.}
An actual point functor $T_a^*(-)\otimes M$ preserving $E$
preserves (N26.1) degree by degree. Distinct component
structure sheaves have no degree-zero morphisms between
them, by (N26.7), and each has only scalar endomorphisms.
The induced cohomology isomorphisms therefore permute the
actual components and have only scalar freedom on each
component. They force $a\in\operatorname{Stab}(Z)$ and
$M|_{C_\nu}$ trivial for every component direction.
On the original $E_i^6$, the coordinate-threefold
directions force $M$ itself trivial, as in HD-N25. The
translation stabilizer has order at most144 by HD-N23.

On a stabilized component, the actual complement addition
isogeny (N25.6) identifies its self-Ext with the exterior
cohomology of $B\times\{d\}$. Translation acts identically
on the invariant Dolbeault factor, and any flat tensor
trivial on $B$ acts identically on the point-Koszul factor.
The finite-etale unit and trace make this calculation
faithful downstairs. Thus every component stabilizer acts
identically on its fifteen-dimensional self-Ext block,
independently of scalar choices of its isomorphism.

For a coherent faithful point group $G$, (N26.3) is an
isomorphism of genuine $G$-modules. Each orbit of actual
components contributes fifteen invariant dimensions.
There are at least $\lceil9408/144\rceil=66$ orbits, and
therefore
\[
                 \dim\operatorname{Ext}^2(E,E)^G\ge990.
                                                               \tag{N26.11}
\]
The full degree-two ambient trace target has dimension495;
the actual invariant trace kernel has dimension at least495.
The identity-sector Hochschild comparison of HD-N12
identifies this with a weak categorical trace kernel for
the equivariant object.

For an actual isogeny $f:X\to A$, flat pullback preserves
the canonical cohomology degrees and replaces each support
by its disjoint connected inverse-image components. Every
such component occurs once, and their cross intersections
remain transverse double points. The entire filtered-Hom
argument applies to $f^*E$. A new preserving point functor
may involve a nontrivial ambient flat line, but its
restriction to each support is trivial; the complement
calculation above still proves identity on a stabilized
self-Ext block. Its translation part maps under $f$ into
$\operatorname{Stab}(Z)$. The actual map from connected
inverse-image components to original component labels is
surjective and equivariant. It induces a surjection on
orbit sets, so there are still at least66 orbits. This
proves (N26.11) and the trace bound for every such new
coherent action, without dividing an Ext bound by the
isogeny degree. Subsequent actual line tensor and
Fourier--Mukai equivalences transport the point group,
Ext maps, and traces by the intertwiners proved in HD-N12
and HD-N17.

\paragraph{5. Verification and scope.}
The original BBD scan was read through Sections 3.1.1--3.1.6,
printed pp.76--79, including the filtered-Hom construction
and proof of the canonical-filtration equivalence. The
source scan, page images, and exact file hashes are saved
in the accompanying source-review directory. The exact
finite index certificate checks the differential positions,
the negative-degree formula on explicit support graphs,
and all integer bounds; it is supplementary to the
universal filtered-map proof above.

The result covers the actual objects with precisely the
cohomology (N26.1), their stated isogeny pullbacks, and
subsequent actual equivalences. It makes no assertion about
cohomology sheaves with repeated copies of a support,
nontrivial bundles on a support, nonreduced support, or
objects whose additional cohomology cancels in the Chern
character. Those require their own actual Ext maps.
